% written by scripts/reuse_short.py from the sources of the short paper; do not edit
The pass rule of \papersec{5} was applied \SGateCount{} to a real candidate: a changed version of the withheld
agent bundle, tested against the unchanged baseline on the same tasks and repeats. Each application ran once. Its rule,
tasks and analysis script were frozen before its first job, and the outcomes stayed unread until the run had ended.
Table~\ref{tab:s-gates} gives every quantity that the rule uses, and Figure~\ref{fig:s-gates} the gain of every
repository.

\paragraph{How the candidates were chosen.}
Each candidate was the leader of a screen on one separate 48-task set, with two repeats per task: 576 runs over the
baseline and five variants for the first candidate, and 288 runs over the baseline and two variants for each of the
other two. Each leader solved more runs than the baseline in its screen (25, 22 and 25 of 96, against 19, 21 and 20).
A screen selects a candidate and tests nothing, because the leader of several variants is expected to look better in
its screen than it is. No task of the screens is among the tasks of an application. Two small pilots came before the
screens of the later candidates. The variants of the second candidate were first tried on 15 tasks of the case-study
panel, which shares no task with an application. Two predecessors of the third candidate were tried on 16 tasks of the
small panel, that is, on tasks of the first application and on none of the second and third.

\paragraph{Tasks and conditions.}
The first application used the small panel (\SGateOneTasks{} tasks) and all five conditions of \papersec{5}. The
second used the \SGateTwoTasks{} tasks of the large panel that are outside the small one. The third used the same
\SGateThreeTasks{} tasks as the second, and both arms ran again. Its registration stated three things before the run:
that this was the second use of that task set, that the candidate had run on none of these tasks, and that the
candidate was built after the result of the second application was known. It also named two mechanical reads of
earlier runs that went into a repair of the candidate and that may include runs on these tasks: a scan of the
agent's commands in the \SGateTwoSessions{} runs of the second application, and the output of one command in six
earlier runs. Neither read used whether a task was solved. The candidate had been submitted to the competition
once before the run; its hosted score had not arrived when the outcomes were opened, and it is not an input of the
rule. The second and third applications registered conditions (1) to (4) only.

\paragraph{Results.}
The rule asks for a mean gain of at least \SGateMinGain{} points with a one-sided $p$ of at most \SGateAlpha, and for
further conditions on repositories, time and infrastructure errors. No candidate passed. The first two estimates are
below zero: \SGateOneGain{} points (\SGateOneCand{} against \SGateOneBase{} solved runs of \SGateOneRuns; $p =
\SGateOneP$) and \SGateTwoGain{} points (\SGateTwoCand{} against \SGateTwoBase{} of \SGateTwoRuns; $p = \SGateTwoP$).
The third candidate solved \SGateThreeMore{} runs more than the baseline (\SGateThreeCand{} against \SGateThreeBase{}
of \SGateThreeRuns), a gain of \SGateThreeGain{} points. This gain is above the required size, and conditions (3) and
(4) are met: \SGateThreeReposUp{} of the \SGateThreeRepos{} repositories gain, \SGateThreeReposDown{} lose, and the
lowest leave-one-repository-out gain is \SGateThreeLoo{} points. But the one-sided $p$ is \SGateThreeP, above
\SGateAlpha, so condition (2) is not met and the decision is FAIL.

\paragraph{How to read the decisions.}
The verdict of a failed application is ``not confirmed''. It is not evidence that the candidate has no effect: the
registration of the third application estimated, before the run, that a run of this size detects a true gain of the
required size about half of the time. Nor do we call the third candidate an improvement. Its $p$ is not corrected for
the three applications, and it is also above \SGateAlphaSplit, the level divided by three. Its gain is uneven: the
three repositories with the largest totals hold \SGateThreeTopThreeTasks{} of the \SGateThreeTasks{} tasks and gain
\SGateThreeTopThree{} task-equivalents, while the net gain over all repositories is \SGateThreeNet{}
(Figure~\ref{fig:s-gates}). A gain concentrated in a few repositories is the case in which the rule has little
power (\papertab{3}), and it is also what a few repositories can produce by chance. The first two leaders did not keep
their screen lead on new tasks; this is consistent with the selection effect that the rule exists to catch.
